\documentclass[11pt,a4paper]{article}

\usepackage[utf8]{inputenc}
\usepackage[T1]{fontenc}
\usepackage{amsmath,amssymb,amsthm}
\usepackage{mathtools}
\usepackage{hyperref}
\usepackage{cleveref}
\usepackage{enumitem}
\usepackage{booktabs}
\usepackage{listings}
\usepackage{xcolor}
\usepackage[margin=1in]{geometry}

% Lean 4 syntax highlighting
\lstdefinelanguage{Lean4}{
  keywords={theorem,def,class,where,extends,inductive,namespace,end,import,
            open,variable,instance,deriving,abbrev,let,have,exact,intro,
            rfl,by,simp,rw,match,with,fun,forall,if,then,else,return},
  keywordstyle=\color{blue}\bfseries,
  ndkeywords={Term,State,Type,Bool,Prop,Nat,Function},
  ndkeywordstyle=\color{teal},
  comment=[l]{--},
  morecomment=[s]{/-}{-/},
  commentstyle=\color{gray}\itshape,
  stringstyle=\color{red},
  sensitive=true,
  morestring=[b]",
  literate={->}{$\to$}2 {!=}{$\neq$}2 {>=}{$\geq$}2
}
\lstset{
  language=Lean4,
  basicstyle=\small\ttfamily,
  breaklines=true,
  frame=single,
  xleftmargin=1em,
  numbers=left,
  numberstyle=\tiny\color{gray},
  captionpos=b,
}

% Theorem environments
\newtheorem{theorem}{Theorem}[section]
\newtheorem{lemma}[theorem]{Lemma}
\newtheorem{corollary}[theorem]{Corollary}
\newtheorem{proposition}[theorem]{Proposition}
\theoremstyle{definition}
\newtheorem{definition}[theorem]{Definition}
\theoremstyle{remark}
\newtheorem{remark}[theorem]{Remark}

% Notation
\newcommand{\compose}{\mathbin{\circ}}
\newcommand{\opl}{\oplus}
\newcommand{\obs}{\mathrm{obs}}
\newcommand{\fzero}{f_0}
\newcommand{\fone}{f_1}
\newcommand{\fn}{f_n}
\newcommand{\Term}{\mathsf{Term}}
\newcommand{\eval}{\mathrm{eval}}
\newcommand{\encode}[1]{\ulcorner #1 \urcorner}
\newcommand{\obsRewrite}{\mathrm{obsRewrite}}
\newcommand{\deepRewrite}{\mathrm{deepRewrite}}
\newcommand{\termSize}{\mathrm{termSize}}
\newcommand{\smul}[2]{#1 \cdot #2}
\newcommand{\iterS}{\mathrm{iterS}}
\newcommand{\N}{\mathbb{N}}
\newcommand{\Z}{\mathbb{Z}}

\title{\textbf{The Encoding Obstruction:}\\ An Impossibility Theorem for Self-Representing Eigenvalue Systems}
\author{[To be determined]}
\date{}

\begin{document}
\maketitle

\begin{abstract}
We prove that no injective encoding of a self-representing algebraic system
can simultaneously preserve eigenvalue arithmetic and commute with the
observation operator. In any commutative semiring with a derivation (the
Leibniz rule), the observation operator produces linear eigenvalues, while
any injective encoding of structured terms requires super-linear pairing
functions. These growth regimes are incompatible: no injective map from
terms to natural numbers is a coalgebra morphism between the semantic
observation (eigenvalue scaling) and the syntactic observation
(rewrite-then-re-encode). We formalize the system as a
\emph{DynSemiring}---a commutative semiring with a derivation, additive
identity, and characteristic zero---and prove all results in Lean~4 with
zero \texttt{sorry}. The impossibility has the structure of a trilemma:
eigenvalue arithmetic, self-encoding, and coalgebra morphism are pairwise
compatible but jointly unsatisfiable, and every escape route (redundant
representations, residue number systems, tropical algebra) sacrifices a
property required for self-encoding. We show that the Leibniz rule is the
unique minimal correct rewrite, that extensionalization more than doubles
term size, and that the intension/extension chain has no fixed point. The
executable consequence is that any self-representing system with addition
must maintain normalization state external to its encoding, because carry
propagation in addition is non-local in any unique positional
representation. We analyze the relationship of these results to the three
known barriers in complexity theory and show that the encoding obstruction
is non-relativizing by construction and plausibly avoids the other two
barriers.
\end{abstract}

\tableofcontents

%% ========================================================================
\section{Introduction}
\label{sec:intro}
%% ========================================================================

\subsection{The Problem}

When a formal system is powerful enough to encode its own syntax as
data---G\"odel encoding in Peano Arithmetic, the universal Turing machine
encoding its own description---a natural question arises: can the system
faithfully simulate its own dynamics through this encoding? That is, does
applying the system's operations to the encoded representation produce the
same result as encoding the result of applying those operations directly?

For \emph{passive} encoding (inert labels assigned to terms), this question
leads to G\"odel's incompleteness theorems: the system can name
self-referential propositions but cannot resolve them. For \emph{executable}
encoding (where the system operates on its own codes using its own
operations), we show the question leads to a stronger impossibility: the
encoding cannot commute with the system's observation operator, and the
failure forces external computational state.

\subsection{The Setting}

We work in a \emph{DynSemiring}: a commutative monoid
$(\mathrm{State}, \compose, \fzero)$ equipped with an observation operator
$\obs$ satisfying the Leibniz (product) rule, an addition operation $\opl$
with additive identity $z$, distributivity of $\compose$ over $\opl$, and
characteristic zero. The complete axiom set is given in
\Cref{def:dynsemiring}.

From these axioms, eigenstates
$\fn = \underbrace{\fone \compose \cdots \compose \fone}_{n}$ satisfy
$\obs(\fn) = \smul{n}{\fn}$, producing the natural numbers as eigenvalues.
This yields a structure equiconsistent with Peano Arithmetic, which provides
the machinery for G\"odel encoding: prime factorization, primitive recursive
predicates, and bounded quantification.

\subsection{The Question}

We define two coalgebra structures on the same carrier:
\begin{itemize}
  \item \textbf{Semantic coalgebra} $(\mathrm{State}, \obs_{\mathrm{sem}})$:
    eigenvalue scaling, $\obs(\fn) = \smul{n}{\fn}$.
  \item \textbf{Syntactic coalgebra} $(\mathrm{State}, \obs_{\mathrm{code}})$:
    decode a G\"odel code to a term, apply the Leibniz rewrite, re-encode.
\end{itemize}

The eval-quote diagram asks whether these produce bisimilar behavior:
\[
\begin{array}{ccc}
  \Term & \xrightarrow{\eval} & \mathrm{State} \\
  \downarrow\scriptstyle{\encode{\cdot}} & & \downarrow\scriptstyle{\obs_{\mathrm{sem}}} \\
  \mathrm{State} & \xrightarrow{\obs_{\mathrm{code}}} & \mathrm{State}
\end{array}
\]
We prove this diagram does not commute. For any injective encoding and any
compound term, the syntactic path produces a different code than the
semantic path.

\subsection{Summary of Results}

All results are formalized in Lean~4 with zero \texttt{sorry}.

\begin{enumerate}[label=\textbf{Theorem \arabic*.}, leftmargin=*]
  \item \emph{(Encoding Obstruction)} For any injective
    $\varepsilon : \Term \to \N$ and any terms $a, b$:
    $\varepsilon(\obsRewrite(\obs(a \compose b)))
    \neq \varepsilon(\obs(a \compose b))$.
  \item \emph{(Semantic Invariance)} For any DynSemiring instance and any
    terms $a, b$:
    $\eval(\obsRewrite(\obs(a \compose b)))
    = \eval(\obs(a \compose b))$.
  \item \emph{(No Fixed Point)} The intensional, one-step extensional, and
    deeply extensional forms are pairwise distinct.
  \item \emph{(Size Growth)} Extensionalization more than doubles term size:
    $\termSize(\obsRewrite(\obs(a \compose b)))
    > 2 \cdot \termSize(\obs(a \compose b))$.
  \item \emph{(Bitwise Impossibility)} No bitwise operation on $\N$ can
    simultaneously satisfy commutativity, cancellation, identity, and
    characteristic zero.
  \item \emph{(Carry Propagation)} For any base $b \geq 2$ and position
    $k \geq 1$, addition produces non-local effects.
\end{enumerate}

%% ========================================================================
\section{The DynSemiring and Term Algebra}
\label{sec:dynsemiring}
%% ========================================================================

\subsection{Definition}

\begin{definition}[DynSemiring]
\label{def:dynsemiring}
A \emph{DynSemiring} is a type $\mathrm{State}$ equipped with operations
$\compose$, $\opl$, $\obs$, and constants $\fzero$, $\fone$, $z$
satisfying:
\begin{enumerate}[label=(DS\arabic*)]
  \item $\fone \compose s = s \compose \fone$ and $\fzero \compose s = s$
    \hfill (monoid)
  \item $z \compose s = z$ \hfill (annihilator)
  \item $a \compose b = b \compose a$ \hfill (commutativity)
  \item $\obs(\fone) = \fone$ \hfill (fixed point)
  \item $\obs(a \compose b) = (\obs(a) \compose b) \opl (a \compose \obs(b))$
    \hfill (Leibniz rule)
  \item $\obs(\fzero) = z$ \hfill (obs of identity)
  \item $\obs(a \opl b) = \obs(a) \opl \obs(b)$ \hfill (obs distributes)
  \item $z \opl a = a$, $a \opl b = b \opl a$,
    $a \opl (b \opl c) = (a \opl b) \opl c$ \hfill (abelian group)
  \item $a \compose (b \opl c) = (a \compose b) \opl (a \compose c)$
    \hfill (distributivity)
  \item $\forall n \geq 1: \underbrace{\fzero \opl \cdots \opl \fzero}_{n} \neq z$
    \hfill (characteristic zero)
\end{enumerate}
\end{definition}

The Lean formalization:

\begin{lstlisting}[caption={DynSemiring definition (Core/Semiring.lean)}]
class DynSemiring (State : Type*) extends SpectralMonoid State where
  oplus : State -> State -> State
  oplus_assoc : forall a b c, oplus a (oplus b c) = oplus (oplus a b) c
  oplus_comm : forall a b, oplus a b = oplus b a
  oplus_z : forall a, oplus z a = a
  compose_distrib : forall a b c, compose a (oplus b c)
                    = oplus (compose a b) (compose a c)
  obs_f0 : obs f0 = z
  leibniz : forall a b, obs (compose a b)
            = oplus (compose (obs a) b) (compose a (obs b))
  obs_oplus : forall a b, obs (oplus a b) = oplus (obs a) (obs b)
  char_zero : forall n : Nat, n >= 1 -> nfoldSum oplus z n f0 != z
\end{lstlisting}

\subsection{The Free Term Algebra}

\begin{lstlisting}[caption={Term inductive type (Obstruction/Encoding.lean)}]
inductive Term where
  | f0 : Term
  | f1 : Term
  | z  : Term
  | compose : Term -> Term -> Term
  | oplus : Term -> Term -> Term
  | obs : Term -> Term
  deriving DecidableEq, Repr
\end{lstlisting}

The evaluation function $\eval : \Term \to \mathrm{State}$ interprets each
constructor in any DynSemiring instance. The key operation is
$\obsRewrite$, which applies the Leibniz rule at the root:

\begin{lstlisting}[caption={One-step extensionalization (Obstruction/Encoding.lean)}]
def obsRewrite : Term -> Term
  | .obs (.compose a b) =>
      .oplus (.compose (.obs a) b) (.compose a (.obs b))
  | .obs (.oplus a b) => .oplus (.obs a) (.obs b)
  | .obs .f1 => .f1
  | .obs .f0 => .z
  | .obs .z  => .z
  | t => t
\end{lstlisting}

\subsection{The Eigenvalue Theorem}

\begin{theorem}[Eigenvalue]
\label{thm:eigenvalue}
For all $n \in \N$:
$\obs(\iterS\; n) = \smul{n}{\iterS\; n}$,
where $\iterS\; n$ denotes the $n$-fold self-composition of $\fone$.
\end{theorem}

\begin{lstlisting}[caption={Eigenvalue theorem (Core/Semiring.lean)}]
theorem eigenvalue (n : Nat) :
    obs (iterS n) = smul n (iterS (State := State) n) := by
  induction n with
  | zero => simp [iterS, obs_f0, smul, nfoldSum]
  | succ n ih =>
    simp only [iterS, succ]
    rw [leibniz, obs_f1, ih, compose_smul]
    rfl
\end{lstlisting}

The eigenvalues $\{0, 1, 2, 3, \ldots\}$ form the natural numbers under
$\opl$ and $\compose$, yielding a structure equiconsistent with PA.

%% ========================================================================
\section{The Encoding Obstruction}
\label{sec:obstruction}
%% ========================================================================

\subsection{The Structural Gap}

\begin{lemma}[Extensionalize $\neq$ Original]
\label{lem:ext-ne-orig}
For any terms $a, b$:
$\obsRewrite(\obs(a \compose b)) \neq \obs(a \compose b)$.
\end{lemma}

\begin{proof}
The left side has constructor $\opl$ at the root (the Leibniz rule produces
a sum). The right side has constructor $\obs$ at the root. These are
distinct constructors in the free term algebra, so the terms are distinct by
$\mathsf{DecidableEq}$.
\end{proof}

\begin{lstlisting}[caption={General structural gap (Obstruction/Encoding.lean)}]
theorem extensionalize_ne_original_general (a b : Term) :
    obsRewrite (.obs (.compose a b)) != .obs (.compose a b) := by
  simp [obsRewrite]
\end{lstlisting}

\begin{theorem}[No Coalgebra Morphism]
\label{thm:no-coalgebra}
For any injective $\varepsilon : \Term \to \N$ and any terms $a, b$:
\[
  \varepsilon(\obsRewrite(\obs(a \compose b)))
  \neq \varepsilon(\obs(a \compose b)).
\]
\end{theorem}

\begin{proof}
By \Cref{lem:ext-ne-orig}, the two terms are distinct. By injectivity of
$\varepsilon$, distinct terms receive distinct codes.
\end{proof}

\begin{lstlisting}[caption={Main theorem (Obstruction/Main.lean)}]
theorem no_coalgebra_morphism
    (e : Term -> Nat) (he : Function.Injective e) (a b : Term) :
    e (obsRewrite (.obs (.compose a b))) != e (.obs (.compose a b)) :=
  encoding_obstruction_general e he a b
\end{lstlisting}

The proof is elementary but the consequence is not: \textbf{no injective
encoding of terms into natural numbers can serve as a coalgebra morphism
between semantic observation and syntactic rewriting.}

\subsection{Semantic Invariance}

\begin{theorem}[Semantic Invariance]
\label{thm:sem-inv}
For any DynSemiring instance and any terms $a, b$:
\[
  \eval(\obsRewrite(\obs(a \compose b))) = \eval(\obs(a \compose b)).
\]
\end{theorem}

\begin{proof}
Direct application of the Leibniz axiom (DS5). The $\obsRewrite$ function
applies this identity at the term level; evaluation in any model respects
the axiom.
\end{proof}

\begin{lstlisting}[caption={Semantic invariance (Obstruction/Encoding.lean)}]
theorem semantic_invariance_general
    {State : Type*} [DynSemiring State] (a b : Term) :
    eval (.obs (.compose a b))
    = (eval (obsRewrite (.obs (.compose a b))) : State) :=
  (obsRewrite_eval (.obs (.compose a b))).symm
-- eval(original) = eval(rewritten) in ALL models
\end{lstlisting}

\Cref{thm:no-coalgebra,thm:sem-inv} together establish the core asymmetry:
\textbf{the encoding changes the syntax while preserving the semantics.}

\subsection{No Fixed Point}

\begin{theorem}[Substitution Chain Distinct]
\label{thm:no-fixpoint}
Let $t_0 = \obs(a \compose b)$,
$t_1 = \obsRewrite(t_0)$, and
$t_2 = \deepRewrite(t_1)$.
Then $t_0$, $t_1$, and $t_2$ are pairwise distinct.
\end{theorem}

\begin{lstlisting}[caption={No fixed point (Obstruction/Encoding.lean)}]
theorem substitution_chain_distinct :
    .obs t_f2 != obsRewrite (.obs t_f2) /\
    obsRewrite (.obs t_f2) != deepRewrite (obsRewrite (.obs t_f2)) /\
    .obs t_f2 != deepRewrite (obsRewrite (.obs t_f2)) :=
  <fun h => absurd h.symm extensionalize_ne_original,
   fun h => absurd h.symm deep_ne_shallow,
   fun h => absurd h.symm deep_ne_original>
\end{lstlisting}

The chain of extensionalizations never stabilizes. Each level of unfolding
produces a genuinely new syntactic form, even though the semantic value is
preserved throughout.

\subsection{Size Growth}

\begin{theorem}[More Than Doubling]
\label{thm:size-growth}
For any terms $a, b$:
\[
  \termSize(\obsRewrite(\obs(a \compose b)))
  > 2 \cdot \termSize(\obs(a \compose b)).
\]
\end{theorem}

\begin{proof}
Let $s_a = \termSize(a)$ and $s_b = \termSize(b)$. Then
$\termSize(\obs(a \compose b)) = 2 + s_a + s_b$ and
$\termSize(\obsRewrite(\obs(a \compose b))) = 5 + 2s_a + 2s_b$.
Since $5 + 2s_a + 2s_b > 2(2 + s_a + s_b) = 4 + 2s_a + 2s_b$, the bound
holds.
\end{proof}

%% ========================================================================
\section{The Trilemma}
\label{sec:trilemma}
%% ========================================================================

\subsection{Three Properties}

Any self-representing eigenvalue system combines three structural properties:

\begin{description}
  \item[Property 1: Eigenvalue Arithmetic.] The Leibniz rule forces
    eigenvalues to be linear: $\obs(\fn) = \smul{n}{\fn}$. This produces
    $\N$, yielding PA-strength arithmetic.
  \item[Property 2: Self-Encoding.] PA arithmetic enables G\"odel encoding:
    an injective, decodable map $\Term \to \N$. Any such encoding requires
    an injective pairing function
    $\pi : \N \times \N \to \N$, which is necessarily super-linear:
    $\pi(n,n) \in \omega(n)$.
  \item[Property 3: Coalgebra Morphism.] The encoding commutes with
    observation---the eval-quote diagram commutes.
\end{description}

\subsection{Pairwise Compatibility, Joint Impossibility}

Properties 1 and 2 are compatible: linear eigenvalues produce arithmetic,
which enables encoding. Properties 2 and 3 would be compatible if
eigenvalues could be made super-linear to match the encoding's growth rate.
But super-linear eigenvalue sets (e.g., $\{1, 2, 4, 8, \ldots\}$) are not
closed under addition, destroying the semiring structure required by
Property~2. Properties 1 and 3 would be compatible if the encoding could be
linear. But a linear function $\N \times \N \to \N$ is not injective on
pairs.

\textbf{The trilemma:} self-encoding requires the arithmetic that only
linear eigenvalues provide, but the coalgebra morphism requires encoding
growth rates to match eigenvalue growth rates, and no injective encoding is
linear.

\subsection{Blocked Escape Routes}

\textbf{Escape 1: Redundant Number Systems.} Signed-digit representations
\cite{avizienis1961} achieve carry-free addition via non-unique
representations. But self-encoding requires injective, decodable maps, which
require unique representation. Carry reappears at equality testing and
canonicalization.

\textbf{Escape 2: Residue Number Systems.} RNS provides carry-free
modular arithmetic, but comparison, ordering, and division are lost.
G\"odel encoding requires positional access, which requires ordering.
Carry reappears at conversion to positional form.

\textbf{Escape 3: Tropical Algebra.} Replacing $(\N, +, \times)$ with
$(\N, \max, +)$ makes addition local, but $\max$ is idempotent and
non-cancellative. Without cancellation: no subtraction, no division, no
prime factorization, no G\"odel encoding.

%% ========================================================================
\section{The Executable Consequence: External State}
\label{sec:external-state}
%% ========================================================================

\subsection{From Obstruction to State}

For passive encoding, the obstruction gives incompleteness. For executable
encoding, it forces external state through the following chain:
\begin{enumerate}
  \item Self-encoding requires addition (the Leibniz rule introduces $\opl$).
  \item Addition on unique positional representations requires carry
    propagation.
  \item Carry propagation is non-local (\Cref{thm:carry}).
  \item Non-local operations require normalization state external to the
    encoding.
\end{enumerate}

\subsection{The Bitwise Impossibility}

\begin{theorem}[Bitwise Impossibility]
\label{thm:bitwise}
No operation $f : \mathrm{Bool} \to \mathrm{Bool} \to \mathrm{Bool}$ can
simultaneously satisfy commutativity, identity, and cancellation without
having characteristic~2.
\end{theorem}

\begin{lstlisting}[caption={Bitwise impossibility (Obstruction/Oplus.lean)}]
theorem oplus_bitwise_impossible :
    forall (f : Bool -> Bool -> Bool) (e : Bool),
      (forall a b, f a b = f b a) ->
      (forall x, f e x = x) ->
      (forall a b c, f a b = f a c -> b = c) ->
      (forall x, f x x = e) :=
  bool_comm_cancel_id_implies_char_two
\end{lstlisting}

Since the DynSemiring requires characteristic zero (axiom DS10),
characteristic~2 is contradictory. No bitwise realization of $\opl$ is
possible.

\subsection{Carry Propagation}

\begin{theorem}[Carry Propagation]
\label{thm:carry}
For any base $b \geq 2$ and position $k \geq 1$, addition can produce
output at position $k$ that depends on inputs at all positions $0$ through
$k-1$.
\end{theorem}

This is formalized in \texttt{Pairing.lean} as
\texttt{carry\_propagation\_general}.

%% ========================================================================
\section{Normalization and Complexity}
\label{sec:normalization}
%% ========================================================================

\subsection{Confluence and Termination}

The DynSemiring axioms, oriented as rewrite rules, form a confluent and
terminating rewriting system. Every term has a unique normal form: a
sum-of-products-of-atoms, the algebraic analogue of expanded polynomial form.

\subsection{Normalization Depth}

For the family
$\mathrm{nestedObsCompose}(d) =
\obs(\mathrm{nestedObsCompose}(d{-}1) \compose \mathrm{nestedObsCompose}(d{-}1))$
with base case $\fone \compose \fone$, normalization depth is
\textbf{constant}---exactly 2 passes of bottom-up rewriting, independent of
nesting depth $d$.

However, term size grows exponentially:
$\termSize(\mathrm{nestedObsCompose}(d)) \geq 2^d$.

The exponential cost lives in per-pass term size, not pass count. The
Leibniz rule's branching doubles the term at each application.

\subsection{The Leibniz Rule Is Forced}

$\obsRewrite$ is the unique minimal correct expansion of
$\obs(a \compose b)$ in any DynSemiring. No smaller term evaluates
correctly in all models: dropping either branch of
$(\obs(a) \compose b) \opl (a \compose \obs(b))$ produces an incorrect
result whenever $\obs(a) \neq z$ and $\obs(b) \neq z$. The branching
factor per rewrite step is algebraically inevitable.

%% ========================================================================
\section{Relationship to Complexity-Theoretic Barriers}
\label{sec:barriers}
%% ========================================================================

\subsection{Relativization}

The encoding obstruction is \textbf{non-relativizing by construction}. The
proof examines internal algebraic structure via the $\obsRewrite$ function,
which pattern-matches on the \texttt{.obs(.compose a b)} constructor. An
oracle query has no constructor; the Leibniz decomposition cannot apply to
opaque terms. This parallels how IP$=$PSPACE is non-relativizing:
arithmetization looks inside formula gates, which fails for oracle gates.

\subsection{Natural Proofs}

The encoding obstruction \textbf{does not operate on truth tables}. The
Razborov--Rudich framework~\cite{razborov1997} concerns properties of
Boolean functions tested on their truth tables. The encoding obstruction
concerns algebraic properties of self-encoding systems: eigenvalue
structure, Leibniz extensionalization, carry propagation. The property is
problem-specific (holds for specific algebraic systems, not a non-negligible
fraction of random functions).

\subsection{Algebrization}

The encoding obstruction uses arithmetic in a structurally different way
from the arithmetization that algebrization~\cite{aaronson2009} captures.
Algebrization concerns proofs that convert Boolean computation into
polynomial evaluation over finite fields and use low-degree extensions. The
encoding obstruction examines growth rates of encoding functions and
structural compatibility of eigenvalue scaling with injective
pairing---properties of the meta-level relationship between encoding and
execution.

The key distinction: algebrization models Boolean-to-arithmetic as a
one-directional extension $\{0,1\} \to \mathbb{F}_q$. In a self-encoding
system, position (arithmetic) and existence (Boolean) are the same layer
viewed at different scales, connected bidirectionally by the derivative. The
algebrization framework does not capture proof techniques exploiting this
bidirectional identity.

\subsection{Assessment}

The encoding obstruction avoids the three known barriers because it examines
the internal algebraic structure of self-encoding systems, rather than
input/output behavior (relativization), truth table properties (natural
proofs), or polynomial evaluation of Boolean functions (algebrization).
Avoiding barriers is necessary but not sufficient for resolving P vs NP.

%% ========================================================================
\section{Discussion}
\label{sec:discussion}
%% ========================================================================

\subsection{What the Results Establish}

The encoding obstruction establishes:
\begin{enumerate}
  \item The eval-quote diagram does not commute for any injective encoding,
    while semantics is preserved. The gap is syntactic, not semantic.
  \item The intension/extension chain has no fixed point.
    Substitution is open-ended.
  \item Each extensionalization step more than doubles term size.
  \item No bitwise realization of addition satisfies the algebraic
    requirements. Characteristic zero is incompatible with Boolean operations.
  \item Carry propagation is non-local in any base.
  \item Every known escape route is blocked.
  \item The Leibniz rule is algebraically forced.
\end{enumerate}

\subsection{What the Results Do Not Establish}

The encoding obstruction does not prove P$\neq$NP. The algebraic separation
(linear vs.\ super-linear) operates at a different threshold than the
complexity-theoretic separation (polynomial vs.\ super-polynomial). The
structural analogy is genuine---verification is intensional, construction is
extensional---but the logical implication is not established. The bridge
from algebraic results to complexity separation requires showing that
normalization cost corresponds to computational complexity for specific
problem encodings. This connection is the central open question.

\subsection{Relationship to G\"odel's Incompleteness}

Both G\"odel's theorem and the encoding obstruction are consequences of the
intension/extension gap in self-naming systems. G\"odel's theorem concerns
passive self-reference: the system names a proposition about itself but
cannot resolve it. The encoding obstruction concerns executable
self-reference: the system attempts to run its own encoding, and the
execution step introduces machinery the encoding cannot absorb.

The relationship is not analogy but specialization. The encoding
obstruction, restricted to passive encoding, gives G\"odel incompleteness.
Extended to executable encoding, it gives the external state requirement.

\subsection{Open Questions}

\begin{enumerate}
  \item \textbf{Algebraic to complexity separation.} Does exponential
    term-size growth under Leibniz extensionalization for NP-complete
    encodings imply super-polynomial time?
  \item \textbf{Weakest behavioral equivalence.} Natural transformation,
    asymptotic equivalence, or partial bisimulation?
  \item \textbf{Generalization.} Do analogous obstructions hold for
    self-representing systems with a homomorphism rule instead of a
    derivation?
  \item \textbf{A formal language for complexity.} Is there a language in
    which contractibility to a ground state encodes complexity class
    membership?
\end{enumerate}

%% ========================================================================
\section{Formalization}
\label{sec:formalization}
%% ========================================================================

All theorems are formalized in Lean~4 using the Mathlib library. The
formalization is organized as follows:

\begin{table}[h]
\centering
\begin{tabular}{ll}
\toprule
\textbf{File} & \textbf{Contents} \\
\midrule
\texttt{Core/Semiring.lean}
  & DynSemiring definition, eigenvalue theorem \\
\texttt{Obstruction/Encoding.lean}
  & obsRewrite, semantic invariance, encoding obstruction, size bounds \\
\texttt{Obstruction/Main.lean}
  & no\_coalgebra\_morphism, trajectory divergence \\
\texttt{Obstruction/Oplus.lean}
  & Bitwise impossibility, characteristic-2 inheritance \\
\texttt{Obstruction/Pairing.lean}
  & Carry propagation, pairing function bounds \\
\texttt{Obstruction/Complexity.lean}
  & Normalization depth, term size bounds \\
\texttt{Obstruction/CostBounds.lean}
  & Per-pass cost analysis, site persistence \\
\texttt{Obstruction/SAT.lean}
  & SAT formula encoding \\
\texttt{Obstruction/VarSharing.lean}
  & Variable sharing analysis \\
\bottomrule
\end{tabular}
\end{table}

The formalization contains zero \texttt{sorry}. Key conjectures disproved by
formalization (e.g., the exponential normalization depth conjecture) have
been corrected and documented.

%% ========================================================================
\bibliographystyle{plain}
\begin{thebibliography}{10}

\bibitem{avizienis1961}
A.~Avizienis.
\newblock Signed-digit number representations for fast parallel arithmetic.
\newblock {\em IRE Trans.\ Electronic Computers}, EC-10(3):389--400, 1961.

\bibitem{aaronson2009}
S.~Aaronson and A.~Wigderson.
\newblock Algebrization: A new barrier in complexity theory.
\newblock {\em ACM Trans.\ Computation Theory}, 1(1):2, 2009.

\bibitem{baker1975}
T.~Baker, J.~Gill, and R.~Solovay.
\newblock Relativizations of the {P=?NP} question.
\newblock {\em SIAM J.\ Computing}, 4(4):431--442, 1975.

\bibitem{goedel1931}
K.~G\"odel.
\newblock \"Uber formal unentscheidbare {S\"atze} der {Principia Mathematica}
  und verwandter {Systeme I}.
\newblock {\em Monatshefte f\"ur Mathematik und Physik}, 38:173--198, 1931.

\bibitem{gurevich2000}
Y.~Gurevich.
\newblock Sequential abstract state machines capture sequential algorithms.
\newblock {\em ACM Trans.\ Computational Logic}, 1(1):77--111, 2000.

\bibitem{hastad1987}
J.~H{\aa}stad.
\newblock {\em Computational Limitations of Small-Depth Circuits}.
\newblock MIT Press, 1987.

\bibitem{mulmuley2001}
K.~Mulmuley and M.~Sohoni.
\newblock Geometric complexity theory {I}.
\newblock {\em SIAM J.\ Computing}, 31(2):496--526, 2001.

\bibitem{razborov1997}
A.~Razborov and S.~Rudich.
\newblock Natural proofs.
\newblock {\em J.\ Computer and System Sciences}, 55(1):24--35, 1997.

\bibitem{turing1936}
A.~M. Turing.
\newblock On computable numbers, with an application to the
  {Entscheidungsproblem}.
\newblock {\em Proc.\ London Math.\ Soc.}, 42:230--265, 1936.

\bibitem{williams2011}
R.~Williams.
\newblock Non-uniform {ACC} circuit lower bounds.
\newblock In {\em IEEE Conf.\ Computational Complexity}, 2011.

\end{thebibliography}

\end{document}
